## 次线性算法 ## 
到这里为止，我非常感兴趣的还是那些有可能获得置信度的算法的性质。我同样确定，要寻找最优预测算法，囤积海量数据是必不可少的。
然而，并不是所有数据都有价值。就像在沙漠深处安装的监视摄像头传来的视频流那样，实际上绝大部分收集来的数据毫无意义。问题在于，在这个大数据的时代，单单读取这些没有意义的数据来验证它们，真的没有意义，所需的时间就已经不切实际了——我就是这样说服自己不去查阅收到的全部邮件，这种做法没问题的!
为了解决这个问题，理论计算机学家将注意力投向了所谓的次线性算法。这些算法的特点就是在不花时间读取所有数据的情况下，仍然能够提取数据集里最核心的内容。也就是说，这些算法能够在读取输入数据时“一目十行”。
这种算法在历史上最典型的例子就是谷歌的算法。在谷歌上搜索几乎瞬间就能得到结果，然而，谷歌的数据库包含了互联网的数百万亿网页中相当大的一部分。对于谷歌来说，在向用户返回结果之前探索整个数据库根本不在考虑范围之内，因为这样做要花上几天时间!因此，谷歌的搜索算法必须是次线性的。
谷歌用到的技巧跟图书馆和词典一样，就是预先对数据库排序和整理，迅速完成对数据库的查阅。比如说，词典将单词按照字母顺序整理，就能让使用者很快知道希望查找的单词应该会在什么地方。更厉害的是，利用字母顺序，给定词典中的一页以及要查找的单词，使用者就能知道这个单词应该在词典中这一页的前面还是后面。一般来说，在(按照全序关系)排序过的数据库中查找数据非常迅速。所谓的二分算法(大概就是你用词典查单词的时候用的方法)的计算时间是数据库大小的对数。
然而谷歌、图书馆管理员和词典使用的算法都需要预先进行大量的计算。要将所有网页整理并排序，除非对其全部访问，否则不可能做到。即使如此，我们能否在既不检索整个数据集，也不预先对其进行处理的前提下，将有用的信息提取出来呢?令人惊讶的是，对于某些有用的信息以及某些数据来说，答案是肯定的。傅里叶变换就属于这种情况。
在19世纪初，约瑟夫·傅里叶等人就曾研究过傅里叶变换，它是一种对声音、图像、视频以及股票走势等信号的描述方式进行变换的方法[8]。我们在这里就只讨论声音信号。
声音可以通过鼓膜附近气压的变化来描述。但还有另一种等价的描述方式，它对于音乐的谱写来说特别简单，那就是利用组成声音的频率来描述这个声音。傅里叶变换就像一个双语词典，能够将声音的振动变化描述转换为频率描述。这样的翻译有着很多用途，一个原因是要改善声音的话，调整频率通常比调整振动更有效，另一个原因就是声音通常只由为数不多的几个频率组成。
实际上，傅里叶变换已经占据了我们的日常生活，据理查德·巴拉纽克教授所言，就连我们的计算机和电话每天都会进行数十亿次傅里叶变换。在每次听音乐、查看数字化图像，或者观看视频的时候，我们都在让机器进行傅里叶变换。
因此，计算傅里叶变换的算法如果有任何加速的可能性，都会在程序计算或者计算所需硬件的方面带来数十亿欧元的收益，更不用说用户节省下来的等待时间了。在1964年，詹姆斯·库利和约翰·图基就完成了一项壮举，大幅缩短了(离散)傅里叶变换的计算时间。尽管最简单的算法需要的时间与数据量的平方成正比，但库利和图基的算法，又叫快速傅里叶变换，可以在接近线性的时间内完成9。也就是说，要对一百万字节(1MB)的数据进行计算的话，快速傅里叶变换只需要数百万次运算(对于现代计算机来说，这相当于几毫秒的计算)，而不是朴素算法所需的100万乘以100万次运算(差不多1分钟的计算)。
9也就是说离散傅里叶变换所需的时间复杂度从O(n^2)变成了O(nlogn)。
然而在大数据的时代，快速傅里叶变换还是太慢了，尤其在处理上十亿字节(GB)甚至上万亿字节(TB)的数据时更是如此。我们能不能让傅里叶变换变得比现在更快?在2012年，哈桑尼耶、因迪克、卡塔比和普赖斯[9]发现，答案是肯定的。也可以说他们提出了一个巧妙的算法，可以对信号最主要的个频率进行近似计算，而需要的时间大概在乘以数据量的对数这个数量级10。重点在于，这个叫作稀疏傅里叶变换的精妙算法的运行时间是次线性的，也就是说它几乎忽略了整个待处理的信号，却能得知这个信号大概是什么。
10它的复杂度实际上是O(klog^2n)。

思考的多种模式
虽然次线性算法在计算机科学中仍不是主流，但我们的大脑似乎喜欢使用这种算法。谁又未曾一目十行读完书面材料，漫不经心地听别人讨论，或者吃饭时根本没有在意食物的味道呢?奇怪的是，有时候文章中的某些词语、讨论中的某些话题或者食物中的某些味道会吸引我们的注意力。当这种情况出现的时候，我们似乎突然就切换到了思考状态。我们会开始使用更缓慢、更精确的算法来仔细分析文章的含义、讨论的深意与本地特色菜的独特味道。
我在这里描述的，正是诺贝尔经济学奖得主丹尼尔·卡内曼的杰作《思考，快与慢》的核心。卡内曼区分了两套思考系统:系统一和系统二。系统一就像稀疏傅里叶变换，非常迅速、有效、勤奋，但有可能出现严重错误;系统二就像精确傅里叶变换，很慢且需要大量能量，懒惰但更正确。
据卡内曼所说，我们通常会将自己等同于系统二，而且几乎不会意识到系统一的存在。然而，在绝大部分时间内主导的都是系统一，而这会让我们经常犯错。试试这个:球拍和球的价格一共是1.1欧元，球拍的价格比球高1欧元，那么球的价格是多少?
很可能你立刻就看到了答案，但这个答案是错的。根据卡内曼所说，如果这种情况出现，那是因为系统一匆匆忙忙给出了第一个想到的答案;只有在之后，也许是当你看到这里的时候，你的系统二才会开始质疑系统一。
我们可能会认为卡内曼希望说服我们放弃系统一，尽可能经常让系统二运作。并非完全如此。毕竟即使系统二得出的结果更正确，但对于大脑来说，使用系统二更疲劳。思考有其代价。然而，正如所有演员、钢琴家和冲浪者都知道的那样，我们不可能一一思考自己的所有行为与动作。我们最终必须能够自然且自发地做出这些动作，无须大量有意识的思考。
要做到无须长久思考就能做出合格的动作，演员、钢琴家和冲浪者用的是同样一套方法:他们的系统二会迫使系统一学习所需的动作。对于那些希望学习或者授课的人来说，这可能是最重要的意见。与其说学习是将信息塞进脑中，不如说是让系统二来教导系统一，使得系统一发现一条捷径，能够迅速解决那些系统二已经知道怎么解决，但要花上许多时间和能量的问题。也就是说，学习的根本就是发现又快又(足够)好的捷径。

==迈进后严谨阶段!==
数学的学习似乎也是如此进行的。在今天，如果有人问我一个关于加法、指数或者图灵机的问题，那么对我来说，不通过系统二就能回答这些问题也并非毫无可能。这是因为，长年以来我的系统二已经成功教会系统一如何在不花费脑力劳动的前提下回答这些问题!我的大脑中储存着大量捷径，能用于解决大量数学问题，即使是那些对于年轻学生来说略感困难的问题。某些人把这种能力称为数学能力，还有些人把它称为数学直觉。我更愿意把它称为系统二通过努力发现的高效捷径。
但这不是系统一的数学训练中最重要的方面。数学家、菲尔兹奖获得者陶哲轩在他的博客中[10]将数学学习分为三个阶段，分别叫作前严谨阶段、严谨阶段和后严谨阶段。简单来说，学生首先会开始摆弄数字和数学概念，而不会忧心于自己执行的那些代数操作的有效性。然后，随着时间流逝以及接受的数学教育越来越多，严谨性的时刻就会到来，而且很快会转变为纯粹主义:一切都应该用形式化的语言做出解释。最后，后严谨阶段属于研究者，他们花上大部分时间构建不同的启发式论证，用来得出定理证明的大体轮廓，其间会暂时舍弃之前学到的形式体系。关键在于，第二个阶段似乎是到达第三个阶段的必经之路。
陶哲轩的说法可以用系统一和系统二的语言来理解。前严谨阶段对应的是系统一和系统二都没有学会严谨性的情况。当系统二发现严谨性及其重要性时，严谨阶段就开始了。在这个阶段中，系统二会教导系统一，让它意识到自身直觉的局限性。但更重要的是，系统一会由此学会测量自身的置信度，这样就能够更好地做出决定，判断自己能解决问题，还是需要向系统二求助。只有经过这种困难的学习过程，系统一才能成为系统二的完美协助者，而不会在重要的时候拖后腿。因此，成为一名优秀的数学家，基本上相当于让系统一足以在这方面胜任——但要达到这一状态，严谨阶段必不可少。
话虽如此，即使到了后严谨阶段，系统二仍会不断教导系统一，让它能够进步。这就是在莱姆基·奥利弗和孙达拉拉詹发表了关于素数最后一位数字的研究结果之后，所有数论学家的大脑中发生的事情。这个发现惊动了数论学家的系统一，它关于素数的直觉基于素数定理，而这一捷径预言相邻素数的最后一位数字大多是独立的。在这个情况
下，这个捷径出错了。自此之后，数论学家可能进行了某种近似贝叶斯推断，以减少素数定理这一捷径在相邻素数的情况中应用的置信度。
丹尼尔·卡内曼的双系统思考模型对于实用贝叶斯主义者来说的确很有趣。当然，这个模型是错的。但毕竟“所有模型都是错的”。然而卡内曼的模型似乎很有用——即使第17章暗示还存在负责创造性思考过程的第三个系统。
在实践中，尤其是在面对大数据的时候，大部分用到的算法大概是迅速的启发式算法，甚至是次线性的。然而拥有更缓慢但更正确的算法仍然是必要的。此外，如果我们拥有数个更缓慢、但我们知道足够正确的算法，那么我们就可以用这些算法来训练那些快速的启发式算法。
假如我能打个赌，(作为合格的贝叶斯主义者，我喜欢打赌!)我会说这就是未来人工智能的模样。